\chapter{牛顿定律}\label{ch:g12:newtons-laws}

悄悄把一台浴室体重秤带进电梯。轿厢向上加速时，
指针读数偏大约三千克；接近顶端时读数偏小；两者之间
的道理很简单——车厢里什么都没变，变的只是运动。本章阐明指针背后的三条定律，
并把它们锻造成贯穿整个力学的方法：
选取系统，画出受力，投影 $\sum \vect F = m\vect a$。

\section{动量与精确化的第一定律}

\begin{definition}[动量]\label{def:g12:newtons-laws:momentum}
质量为 $m$（\unit{kg}）、速度为 $\vect v$ 运动的物体，其\emph{动量}\index{动量}
是矢量
\[
\vect p = m \vect v \qquad (\unit{kg.m/s}):
\]
即以质量加权的速度；对于系统，则为各部分动量的矢量和。
\end{definition}

\begin{definition}[惯性参考系]\label{def:g12:newtons-laws:inertial}
\emph{惯性参考系}\index{惯性参考系}是这样一种参考系：其中每个孤立物体——不受力，或所受力相互抵消——
都沿直线做匀速运动（见\cref{ch:g10:inertia}）。地面在相当高的精度上是惯性系，
沿直线匀速运动的任何车厢在地面上看也是惯性系；
刹车中的火车或转动的旋转木马则不是：
此时松散的行李无人推动却会自行加速。
\end{definition}

\begin{theorem}[牛顿第一定律]\label{thm:g12:newtons-laws:first}\index{牛顿第一定律}
在惯性系中，孤立系统的动量保持不变：$\vect p$
为常量，即 $d\vect p/dt = \vect 0$。
\end{theorem}
\admitted

\begin{remark}[惯性原理的两次精确化]\label{rem:g12:newtons-laws:sharpened}
\cref{ch:g10:inertia}中的惯性原理经过两次提升：它说明了
\emph{何处}成立——它把惯性系定义为惯性原理成立的那些参考系——
并且它约束的是矢量 $m\vect v$，这一记账方式在本章末尾显示出价值。
\end{remark}

\section{牛顿第二定律}

\begin{theorem}[牛顿第二定律]\label{thm:g12:newtons-laws:second}\index{牛顿第二定律}
在惯性系中，物体所受合力等于其
动量的变化率；对于质量恒定的情形，由 $\vect p = m\vect v$ 可得
\[
\sum \vect F = \frac{d\vect p}{dt} = m\,\frac{d\vect v}{dt} = m \vect a .
\]
\end{theorem}
\admitted

\begin{remark}[自然定律及其单位]\label{rem:g12:newtons-laws:nature}
这一定律无法证明：它是一条自然定律，三个世纪以来一直接受
实验的检验。去年（见\cref{ch:g11:forces-and-motion}）给出了它的影子，
$\Delta v \propto F\,\Delta t/m$；导数（见\cref{ch:g12:kinematics-2d}）使其
在每一时刻都精确成立——其细节留待本科一年级分册。力有了单位，
$\unit{N} = \unit{kg.m/s^2}$；仅受重力作用的物体满足
$m\vect a = m\vect g$：任何质量的物体都以 $\vect a = \vect g$ 下落，且
$g = \qty{9.81}{N/kg} = \qty{9.81}{m/s^2}$——两种读法是一致的。
\end{remark}

\begin{example}[平均力，无需细节]\label{ex:g12:newtons-laws:average}
一辆 \qty{1300}{kg} 的汽车由静止加速到 \qty{100}{km/h}
（\qty{27.8}{m/s}），用时 \qty{8.0}{s}，无论逐时刻的过程如何，合力的平均值为
$\Delta p/\Delta t = 1300 \times 27.8/8.0 \approx \qty{4.5e3}{N}$——动量把一个力的问题
化为前后两个状态的记账。
\end{example}

\section{牛顿第三定律}

\begin{theorem}[牛顿第三定律]\label{thm:g12:newtons-laws:third}\index{牛顿第三定律}
若物体 $A$ 施加一个力 $\vect F_{A \to B}$ 作用于物体 $B$，则 $B$ 对
$A$ 施加的力为 $\vect F_{B \to A} = -\vect F_{A \to B}$：作用线相同，
大小相等，方向相反——在每一时刻都成立，无论运动状态如何，
无论是接触作用还是超距作用。
\end{theorem}
\admitted

\begin{example}[书、桌子与地球]\label{ex:g12:newtons-laws:book}
书放在桌上：两组相互作用，两对力。引力方面：地球把书
向下拉（$\vect P$），书把地球向上拉。接触方面：桌子把书向上推
（$\vect N$），书把桌子向下压。行走的道理相同：
向后蹬地；地面的反作用力是唯一推动你向前的力。
\end{example}

\begin{omfigure}
\begin{tikzpicture}[scale=0.9, font=\small]
  \draw[thick] (0,0) -- (4.6,0); \draw (0.5,0) -- (0.5,-1.4) (4.1,0) -- (4.1,-1.4);
  \draw[fill=omExo!15] (1.6,0) rectangle (2.6,0.6);
  \draw[-Stealth, omThm, very thick] (1.9,0.3) -- ++(0,-1.7) node[below] {$\vect P$};
  \draw[-Stealth, omProp, very thick] (2.3,0.3) -- ++(0,1.7) node[above] {$\vect N$};
  \draw[-Stealth, omProp, very thick] (3.4,-0.05) -- ++(0,-1.7) node[below] {$-\vect N$};
  \draw[thick] (0,-3.1) -- (4.6,-3.1);
  \foreach \x in {0.4,0.8,...,4.5} \draw[gray] (\x,-3.1) -- (\x-0.25,-3.4);
  \node[above right] at (3.9,-3.05) {地球};
  \draw[-Stealth, omThm, very thick] (0.9,-3.1) -- ++(0,1.7) node[above] {$-\vect P$};
\end{tikzpicture}

{\small 一本书，两组相互作用，共四个箭头：引力对（地球拉书，
书拉地球）与接触对（桌子推书，书压桌子）。
$\vect P$ 与 $\vect N$ \emph{并非}一对作用与反作用力：
它们作用于同一本书，且只有在书的加速度为零时才会相互抵消。}
\end{omfigure}

\section{方法：选取、画图、投影}

\begin{method}[求解动力学问题]\label{met:g12:newtons-laws:method}
\begin{enumerate}
  \item \emph{系统}：指明所研究的物体（或物体组）。
  \item \emph{参考系}：选取惯性系，并选择与运动相适应的坐标轴——
        斜面坐标轴，或曲线运动所用的弗勒内标架（见\cref{ch:g12:kinematics-2d}）。
  \item \emph{列举受力}：列出每一个力——先是重力，再是每个接触面上的一个力，
        除此之外别无其他——并画出受力图：
        一个点，每个力一根箭头。
  \item \emph{投影}：写出 $\sum \vect F = m\vect a$，并将其投影到各坐标轴上。
  \item \emph{求解与检验}：检查单位、符号，以及极限情形（$\mu \to 0$，
        $\alpha \to 0$，\dots）。
\end{enumerate}
\end{method}

\begin{definition}[法向力与摩擦力]\label{def:g12:newtons-laws:friction}
一个表面通过两个分量作用于物体：\emph{法向力}\index{法向力
} $\vect N$ 垂直于表面，摩擦力 $\vect f$ 与表面相切。
实验给出摩擦定律：物体未滑动时，静摩擦力自行调节，
直至达到上限，$f_s \leq \mu_s N$；物体一旦滑动，动摩擦力
为 $f_k = \mu_k N$，方向与滑动方向相反。\emph{摩擦因数}%
\index{摩擦因数} $\mu_s$ 与 $\mu_k$（$\mu_k \leq \mu_s$）
仅取决于两种材料。
\end{definition}

\begin{example}[斜面问题两次求解]\label{ex:g12:newtons-laws:incline}
一个物块沿倾角为 $\alpha$ 的斜面滑下；坐标轴：$x$ 沿斜面向下，
$y$ 轴垂直于斜面。对 $y$ 方向：$N = mg\cos\alpha$。不计摩擦时，对 $x$ 方向：$ma = mg\sin\alpha$，
于是 $a = g\sin\alpha$——与质量无关；当 $\alpha = 30^\circ$ 时，$a = \qty{4.9}{m/s^2}$。
若计入动摩擦力，则 $a = g(\sin\alpha - \mu_k\cos\alpha)$；取 $\mu_k = 0.20$，
$a = 9.81\,(0.500 - 0.173) = \qty{3.2}{m/s^2}$。物块原本之所以能静止，
只有在所需的静摩擦力不超过其上限时才成立：$\tan\alpha \leq \mu_s$。
\end{example}

\begin{omfigure}
\begin{tikzpicture}[scale=1.0, font=\small]
  \draw[gray!60] (-0.3,0) -- (6.4,0); \draw[thick] (0,0) -- (30:6.2);
  \draw (1.6,0) arc (0:30:1.6); \node at (14:2.0) {$\alpha$};
  \begin{scope}[shift={(30:3.6)}, rotate=30]
    \draw[fill=omExo!15] (-0.6,0) rectangle (0.6,0.75);
    \coordinate (G) at (0,0.375); \fill (G) circle (1.5pt);
    \draw[-Stealth, omDef, dashed, thick] (G) -- ++(-2.3,0) node[above] {$x$};
    \draw[-Stealth, omDef, dashed, thick] (G) -- ++(0,1.9) node[right] {$y$};
    \draw[-Stealth, omThm, very thick] (G) -- ++(0,1.5) node[left] {$\vect N$};
    \draw[-Stealth, omThm, very thick] (G) -- ++(1.35,0) node[above] {$\vect f$};
  \end{scope}
  \draw[-Stealth, omThm, very thick] (G) -- ++(0,-1.9) node[below] {$\vect P$};
  \draw[-Stealth, omExo, dashed] (G) -- ++(210:0.95) (G) -- ++(-60:1.645);
\end{tikzpicture}

{\small 斜面上的受力图：重力 $\vect P$、法向力 $\vect N$，
摩擦力 $\vect f$ 沿斜面向上（物块下滑时）。沿斜面与垂直斜面的
坐标轴将 $\vect P$ 分解为 $mg\sin\alpha$ 与 $mg\cos\alpha$（虚线）。}
\end{omfigure}

\begin{definition}[视重]\label{def:g12:newtons-laws:apparent}
物体的\emph{视重}\index{视重}是物体与其支持物之间的
法向力——即其下方秤的示数。
\end{definition}

\begin{example}[电梯]\label{ex:g12:newtons-laws:elevator}
一名 \qty{70}{kg} 的乘客站在电梯中的秤上，电梯竖直方向的加速度为
$a_z$（取向上为正）：$N - mg = ma_z$，于是 $N = m(g + a_z)$。向上加速时，
$a_z = \qty{1.5}{m/s^2}$，示数 \qty{792}{N}，显得较重；接近顶层减速时，示数 \qty{582}{N}，显得较轻；
匀速运行时：$mg = \qty{687}{N}$——此即牛顿第一定律；自由落体时，$a_z = -g$：
$N = 0$，秤上读数为零，即完全失重。
\end{example}

\begin{example}[两个物块与一个滑轮]\label{ex:g12:newtons-laws:pulley}
放在光滑桌面上的一个物块 $m_1 = \qty{4.0}{kg}$，通过一根轻绳，跨过
一个理想滑轮，与悬挂的物块 $m_2 = \qty{1.0}{kg}$ 相连。绳子在两端传递
相同的张力 $T$；两物块的加速度大小相同，均为 $a$。
对每个物块沿其运动方向分别应用牛顿第二定律：$m_1 a = T$ 与 $m_2 a = m_2 g - T$，于是
$a = m_2\,g/(m_1 + m_2) = \qty{1.96}{m/s^2}$，并且 $T = m_1 a = \qty{7.8}{N}$——
小于 $m_2 g = \qty{9.8}{N}$，这正是 $m_2$ 向下加速所必需的。
\end{example}

\begin{omfigure}
\begin{tikzpicture}[scale=0.95, font=\small]
  \draw[thick] (0,1.1) -- (4.0,1.1); \draw (0.4,1.1) -- (0.4,-0.4) (3.6,1.1) -- (3.6,-0.4);
  \draw[fill=omExo!15] (1.0,1.1) rectangle (2.0,1.75); \node at (1.5,1.42) {$m_1$};
  \draw (4.0,1.1) -- (4.25,1.35);
  \draw[fill=white] (4.25,1.35) circle (0.25); \fill (4.25,1.35) circle (1.2pt);
  \draw[thick] (2.0,1.6) -- (4.25,1.6) (4.5,1.35) -- (4.5,0.15);
  \draw[fill=omExo!15] (4.2,-0.45) rectangle (4.8,0.15); \node at (4.5,-0.15) {$m_2$};
  \draw[-Stealth, omThm, very thick] (2.0,1.42) -- ++(1.2,0) node[pos=0.8, below] {$\vect T$};
  \draw[-Stealth, omThm, very thick] (4.5,0.15) -- ++(0,0.85) node[right] {$\vect T$};
  \draw[-Stealth, omThm, very thick] (4.5,-0.15) -- ++(0,-1.1) node[right] {$\vect P_2$};
  \draw[-Stealth, omDef, thick] (1.1,2.15) -- ++(0.8,0) node[right] {$\vect a$};
  \draw[-Stealth, omDef, thick] (5.4,0.4) -- ++(0,-0.8) node[below] {$\vect a$};
\end{tikzpicture}

{\small 两个物块、一根轻绳、一个理想滑轮：相同的张力 $T$ 同时拉住两端，
两物块加速度大小相同——每个物块应用一次牛顿第二定律，得到两个
方程，两个未知数。}
\end{omfigure}

\begin{example}[水平弯道]\label{ex:g12:newtons-laws:curve}
质量为 $m$ 的汽车驶过半径为 $R$ 的水平弯道，速率恒为 $v$。在
弗勒内标架（见\cref{ch:g12:kinematics-2d}）中，加速度纯为法向，
$a_N = v^2/R$，指向圆心。竖直方向上，$N = mg$；水平方向上，唯一可用的
向心力是路面作用于轮胎的静摩擦力：
$f = mv^2/R \leq \mu_s N = \mu_s mg$，于是 $v \leq \sqrt{\mu_s\, g R}$，
与质量无关。取 $R = \qty{90}{m}$，干燥路面（$\mu_s = 0.70$）：
$v_{\max} \approx \qty{25}{m/s}$（\qty{89}{km/h}）；在冰面上，上限骤降——
见\cref{exo:g12:newtons-laws:13}，路面做成倾斜的；见\cref{exo:g12:newtons-laws:10}，
同一投影思路用于绳上的摆动。
\end{example}

\begin{omfigure}
\begin{tikzpicture}[scale=0.85, font=\small]
  \draw[gray!50, dashed] (0,0) circle (2.1); \fill (0,0) circle (1.5pt);
  \draw[dashed, omExo] (0,0) -- (150:2.1) node[midway, above] {$R$};
  \draw[fill=omExo!15, rotate around={90:(2.1,0)}] (1.6,-0.28) rectangle (2.6,0.28);
  \draw[-Stealth, omDef, thick] (2.1,0.55) -- ++(0,1.4) node[above] {$\vect v$};
  \draw[-Stealth, omThm, very thick] (2.38,0) -- ++(-1.45,0) node[pos=0.85, above] {$\vect f$};
  \node at (0.2,-2.7) {俯视图};
\end{tikzpicture}
\qquad
\begin{tikzpicture}[scale=0.85, font=\small]
  \draw[thick] (-0.5,0) -- (3.7,0); \draw[fill=omExo!15] (0.9,0.42) rectangle (2.5,1.25);
  \draw[fill=white] (1.2,0.42) circle (0.26) (2.2,0.42) circle (0.26);
  \coordinate (C) at (1.7,0.85); \fill (C) circle (1.2pt);
  \draw[-Stealth, omThm, very thick] (C) -- ++(0,1.4) node[above] {$\vect N$};
  \draw[-Stealth, omThm, very thick] (C) -- ++(0,-1.7) node[below] {$\vect P$};
  \draw[-Stealth, omThm, very thick] (C) -- ++(-1.6,0) node[above] {$\vect f$};
  \node at (1.6,-2.7) {后视图（中心在左侧）};
\end{tikzpicture}

{\small 汽车以恒定速率通过平直弯道：从上方看（左图），速度
沿切线方向，合力指向圆心；从后方看（右图），$\vect N$ 与
$\vect P$ 平衡，路面对车轮的侧向静摩擦力提供全部
向心力。}
\end{omfigure}

\section{孤立系统：动量守恒}

\begin{proposition}[动量守恒]\label{prop:g12:newtons-laws:conservation}\index{动量守恒}
若两个物体所受外力为零或相互抵消，则总动量
$\vect p_1 + \vect p_2$ 恒定不变，与它们彼此施加的力无关。
\end{proposition}

\begin{proof}
将两物体的牛顿第二定律相加：$\frac{d}{dt}(\vect p_1 + \vect p_2) = \vect F_{2 \to 1} +
\vect F_{1 \to 2} + \vect F_{\text{外}}$：一对内力依第三定律相互抵消，
由假设 $\vect F_{\text{外}} = \vect 0$。
\end{proof}

\begin{example}[反冲]\label{ex:g12:newtons-laws:recoil}
一名 \qty{60}{kg} 的滑冰者静止在光滑冰面上，抛出一个 \qty{4.0}{kg} 的药球
以 \qty{6.0}{m/s} 的速率水平抛出。抛出前：$\vect p = \vect 0$；抛出后：
$0 = 60\,V + 4.0 \times 6.0$，于是 $V = \qty{-0.40}{m/s}$——她向后滑去。
没有外力推动她：推她的是药球（第三定律），系统的总动量仍为零。
\end{example}

\begin{remark}[火箭如何不借助外物推进]\label{rem:g12:newtons-laws:rocket}
火箭是一台抛射质量的机器。每秒，发动机向后喷出一批气体，
气体带走向后的动量，火箭便获得等量的
向前动量——即稳定的推力，其大小约为每秒喷出的质量乘以
喷射速率。火箭推的是自身喷出的气体，而非空气；
因此在真空中效果最好。
\end{remark}

\section{练习}

\begin{exercise}[$\star$]\label{exo:g12:newtons-laws:1}
计算下列物体的动量：(a) 一名 \qty{70}{kg} 的短跑运动员以 \qty{10}{m/s} 运动；(b)
一辆 \qty{1300}{kg} 的汽车以 \qty{50}{km/h} 行驶；(c) 一颗 \qty{8.0}{g} 的子弹以 \qty{800}{m/s} 飞行。
将它们按动量大小排序。为什么速度最快的排在最后？
\end{exercise}

\begin{exercise}[$\star$]\label{exo:g12:newtons-laws:2}
一名 \qty{70}{kg} 的乘客站在电梯内的台秤上。当轿厢
(a) 以 \qty{1.2}{m/s^2} 的加速度向上加速，(b) 以 \qty{2.0}{m/s}
的速度向上匀速运动，(c) 以 \qty{1.2}{m/s^2} 的加速度向下加速时，台秤示数各为多少？哪一条定律可以确定 (b) 的结果？
\end{exercise}

\begin{exercise}[$\star$]\label{exo:g12:newtons-laws:3}
一物块自倾角为 $30^\circ$ 的光滑斜面上被释放。计算其
加速度，再计算其沿斜面滑行 \qty{3.0}{m} 后的速率。质量去哪儿了？
\end{exercise}

\begin{exercise}[$\star$]\label{exo:g12:newtons-laws:4}
一个 \qty{58}{g} 的网球以 \qty{15}{m/s} 的速率飞来，被沿同一直线
以 \qty{25}{m/s} 的速率击回；接触时间为 \qty{5.0}{ms}。计算动量的变化量
以及拍线对球的平均作用力，并将该力与球的重力相比较。
\end{exercise}

\begin{exercise}[$\star$]\label{exo:g12:newtons-laws:5}
一本书放在桌面上。写出作用于书上的两个力，并对每个力指出其
作用–反作用力中的对应力（哪个物体？方向如何？）。为什么书所受的重力
与桌面对书的支持力不互为对应力？
\end{exercise}

\begin{exercise}[$\star\star$]\label{exo:g12:newtons-laws:6}
一个 \qty{50}{kg} 的箱子放在水平地面上，$\mu_s = 0.45$，$\mu_k = 0.35$。
(a) 多大的水平推力恰好使其开始运动？(b) 若保持该推力，其
加速度为多少？(c) 多大的推力使其保持匀速运动？
\end{exercise}

\begin{exercise}[$\star\star$]\label{exo:g12:newtons-laws:7}
同一箱子放在可调斜面上（$\mu_s = 0.45$，$\mu_k = 0.35$）。(a) 倾角
多大时，箱子开始下滑？(b) 倾角为 $30^\circ$ 时，计算其加速度。
(c) 若在 (a) 所求的倾角下箱子已开始下滑，它是否会持续加速？为什么？
\end{exercise}

\begin{exercise}[$\star\star$]\label{exo:g12:newtons-laws:8}
一个 \qty{3.0}{kg} 的物块放在光滑桌面上，通过理想滑轮与一
个悬挂的 \qty{2.0}{kg} 物块相连。计算加速度与张力。为什么
张力必然小于悬挂物块的重力？
\end{exercise}

\begin{exercise}[$\star\star$]\label{exo:g12:newtons-laws:9}
一段半径为 \qty{50}{m} 的平直弯道。计算汽车在该弯道上能达到的最大速率，
取 $\mu_s = 0.70$（干燥路面）与 $0.40$（潮湿路面）时，最大速率以 \unit{km/h} 为单位。为何该结果对
轻便摩托车与卡车都成立？
\end{exercise}

\begin{exercise}[$\star\star$]\label{exo:g12:newtons-laws:10}
一个 \qty{0.30}{kg} 的摆球系于 \qty{1.2}{m} 长的细线上，在水平面内
做圆周运动，细线始终与竖直方向保持 $25^\circ$ 的夹角（圆锥摆）。将
第二定律先沿竖直方向、再沿半径方向投影，计算张力与速率。
\end{exercise}

\begin{exercise}[$\star\star$]\label{exo:g12:newtons-laws:11}
一个 \qty{5.0}{kg} 的包裹放在货梯中的台秤上，读数如下：
\begin{center}\small\begin{tabular}{lccc}
阶段 & \qtyrange{0}{2}{s} & \qtyrange{2}{8}{s} & \qtyrange{8}{10}{s}\\
台秤示数（\unit{N}） & 54.0 & 49.0 & 44.1
\end{tabular}\end{center}
计算各阶段的加速度，并描述一种可能的运动过程。电梯
整个过程中是否可能一直\emph{向下}运动？
\end{exercise}

\begin{exercise}[$\star\star\star$]\label{exo:g12:newtons-laws:12}
一名 \qty{60}{kg} 的滑冰者静止在光滑冰面上，将一个 \qty{4.0}{kg} 的球
沿水平方向抛出。(a) 球\emph{相对地面}以 \qty{8.0}{m/s} 的速率飞出时，她的
反冲速率为多少？(b) 若 \qty{8.0}{m/s} 是\emph{相对她}测得的，
则重新计算两个速率。(c) 是哪一个力使她向后运动？
\end{exercise}

\begin{exercise}[$\star\star\star$]\label{exo:g12:newtons-laws:13}
一段公路弯道半径为 \qty{80}{m}，设计车速为 \qty{20}{m/s}，
且不依靠摩擦力。证明路面须以角度 $\theta$ 倾斜，
该角度满足 $\tan\theta = v^2/(gR)$，并算出其值。若汽车以较慢——或较快的速度通过弯道，
摩擦力应起什么作用？
\end{exercise}

\begin{exercise}[$\star\star\star$]\label{exo:g12:newtons-laws:14}
火箭发动机每秒以 \qty{2500}{m/s} 的速度喷出 \qty{250}{kg} 气体。(a) 每秒
带走多少动量——燃气反过来施加的推力为多大？(b) 火箭点火时质量为
\qty{5.0e4}{kg}：其重力为多大？其离地瞬间的加速度为多大？(c) 为什么
发动机在真空中同样能工作？
\end{exercise}

\begin{exercise}[$\star\star\star$]\label{exo:g12:newtons-laws:15}
桌面上质量为 \qty{4.0}{kg} 的物块（$\mu_s = 0.35$，$\mu_k = 0.25$），
通过理想滑轮与一悬挂的质量 $m_2$ 相连。(a) 至少需要多大的 $m_2$ 才能使物块开始运动？
(b) 当 $m_2 = \qty{2.0}{kg}$ 时，计算加速度与张力。(c) 验证当 $\mu_k = 0$ 时，
结果应回到\cref{exo:g12:newtons-laws:8}中的无摩擦公式。
\end{exercise}

\section{问题：缆车}

\begin{problem}[{周末问题——缆车认证：静止时斜坡所需的张力、启动时的附加张力，
越过塔顶时的减载、紧急制动的代价，
以及车厢内标牌可承诺的那一个数字}]
\label{pb:g12:newtons-laws:1}
一辆山地缆车——车厢连同车架，空载质量 $M_0 = \qty{2.0e3}{kg}$——
沿倾角为 $\alpha = 30^\circ$ 的承载索滚动，由与斜坡平行的牵引索
以张力 $\vect T$ 拉动。滚动摩擦可忽略，车厢地板全程保持水平。
车厢可容纳 $25$ 名平均质量为
\qty{80}{kg} 的乘客；你是认证工程师。除非另有说明，车厢
按满载计，即 $M = \qty{4.0e3}{kg}$。

\medskip
\noindent\textbf{第一部分——静止在斜坡上。}
\begin{enumerate}
  \item 选取研究对象；说明为何山体参考系可以视为惯性参考系；
        列出三个力并注明方向。
  \item 静止等待时，遵循哪一条定律？将其沿斜坡方向及垂直斜坡方向投影，
        以表示 $T$ 及承载索对车厢的法向力 $N$。
  \item 计算 $T$ 和 $N$。
  \item 重新计算空车时的 $T$。为什么 $T$ 与总质量成正比？
  \item 牵引索断裂、制动解除，证明失控加速度为
        $g \sin\alpha$（无论满载或空载），并计算其值。
\end{enumerate}

\medskip
\noindent\textbf{第二部分——启动。}卷扬机使满载车厢以
$a = \qty{0.60}{m/s^2}$ 沿斜坡向上加速，直至以 $v = \qty{6.0}{m/s}$ 匀速行驶。
\begin{enumerate}[resume]
  \item 表示并计算此阶段的张力 $T'$。
  \item 启动时张力比静止时增加了多少（用比例表示）？
  \item 此阶段持续多长时间？经过多大距离？
  \item 一名质量为 \qty{80}{kg} 的乘客站在水平地板上。将牛顿第二定律竖直投影，
        以表示并计算她的视重，并与 $mg$ 比较。
  \item 哪一个力使她水平加速？大小为多少？乘客为何会倾斜身体？
\end{enumerate}

\medskip
\noindent\textbf{第三部分——越过塔顶。}承载索以恒定的 \qty{6.0}{m/s} 速率运动，
沿半径为 $R = \qty{40}{m}$ 的竖直圆弧越过支承塔；
在塔顶处，速度沿水平方向。
\begin{enumerate}[resume]
  \item 在塔顶处，取\cref{ch:g12:kinematics-2d}中的弗勒内标架，给出加速度的两个分量，
        并给出它们的数值。
  \item 将第二定律沿法线方向投影，以表示并计算
        承载索对车架的法向力 $N'$。
  \item 同样地，求塔顶处乘客的视重：她的视重
        相对其重力减轻了多少比例？
  \item 速率为多大时 $N'$ 将为零？与匀速运行速率比较，并加以评论。
  \item 车架对越过塔顶的承载索施加了什么力——依据哪条定律，
        其大小与方向各为多少？
\end{enumerate}

\medskip
\noindent\textbf{第四部分——制动、跳跃与标牌。}
\begin{enumerate}[resume]
  \item 在水平的到站路段，紧急制动使满载车厢
        在 \qty{1.5}{s} 内由 \qty{6.0}{m/s} 减速至静止：将 $\sum \vect F = d\vect p/dt$
        理解为制动过程中的平均值，以计算制动力。
  \item 要使一名 \qty{80}{kg} 的乘客停下，需要多大的水平力？
        鞋底–地板间的静摩擦（$\mu_s = 0.40$）能否提供该力？请作出结论。
  \item 在站台处，制动解除，空车悬挂着静止不动；最后一名
        乘客（\qty{80}{kg}）以 \qty{2.5}{m/s} 的速度沿水平方向跳下，
        \emph{相对车厢}的速度。计算车厢的反冲速率与
        乘客对地面的速率。
  \item 用一句话回答：哪一种推进装置依靠这一反冲原理工作，
        其推力由什么决定？
  \item \emph{认证。}牵引索在 \qty{100}{kN} 时断裂，规范将
        工作张力限制为其五分之一。证明上文遇到的最不利情形为
        启动阶段，由此推出允许的最大总质量，并写出标牌内容：
        可载多少名 \qty{80}{kg} 的乘客？并校核安全系数。
\end{enumerate}
\end{problem}
